% ==== 标题页 ====
\begin{titlepage}
\thispagestyle{front}\pagenumbering{gobble}
\vspace*{30mm}
\begin{center}
{\Large K\^osaku Yosida\par}
\vspace{12mm}
{\huge\bfseries Functional Analysis\par}
\vspace{10mm}
{\large Sixth Edition\par}
\vspace{14mm}
{\large Springer-Verlag\par Berlin Heidelberg New York\par}
{\large 1980\par}
\end{center}
\end{titlepage}
% ==== 序言 ====
\pagenumbering{roman}\setcounter{page}{5}
\vspace{6mm}\begin{center}\bfseries Preface to the First Edition\end{center}\vspace{2mm}
The present book is based on lectures given by the author at the University of Tokyo during the past ten years. It is intended as a textbook to be studied by students on their own or to be used in a course on Functional Analysis, i.e., the general theory of linear operators in function spaces together with salient features of its application to diverse fields of modern and classical analysis. 

Necessary prerequisites for the reading of this book are summarized, with or without proof, in Chapter~\hyperref[chap:0]{$0$} under titles: Set Theory, Topological Spaces, Measure Spaces and Linear Spaces. Then, starting with the chapter on Semi-norms, a general theory of Banach and Hilbert spaces is presented in connection with the theory of generalized functions of S. L. \textsc{Sobolev} and L. \textsc{Schwartz}. While the book is primarily addressed to graduate students, it is hoped it might prove useful to research mathematicians, both pure and applied. The reader may pass, e.g., from Chapter~\hyperref[chap:IX]{IX} (Analytical Theory of Semi-groups) directly to Chapter~\hyperref[chap:XIII]{XIII} (Ergodic Theory and Diffusion Theory) and to Chapter~\hyperref[chap:XIV]{XIV} (Integration of the Equation of Evolution). Such materials as ``Weak Topologies and Duality in Locally Convex Spaces'' and ``Nuclear Spaces'' are presented in the form of the appendices to Chapter~\hyperref[chap:V]{V} and Chapter~\hyperref[chap:X]{X}, respectively. These might be skipped for the first reading by those who are interested rather in the application of linear operators. 

In the preparation of the present book, the author has received valuable advice and criticism from many friends. Especially, Mrs. K. \textsc{Hille} has kindly read through the manuscript as well as the galley and page proofs. Without her painstaking help, this book could not have been printed in the present style in the language which was not spoken to the author in the cradle. The author owes very much to his old friends, Professor E. \textsc{Hille} and Professor S. \textsc{Kakutani} of Yale University and Professor R. S. \textsc{Phillips} of Stanford University for the chance to stay in their universities in $1962$, which enabled him to polish the greater part of the manuscript of this book, availing himself of their valuable advice. Professor S. \textsc{Ito} and Dr. H. \textsc{Komatsu} of the University of Tokyo kindly assisted the author in reading various parts of the galley proof, correcting errors and improving the presentation. To all of them, the author expresses his warmest gratitude. 

Thanks are also due to Professor F. K. \textsc{Schmidt} of Heidelberg University and to Professor T. \textsc{Kato} of the University of California at Berkeley who constantly encouraged the author to write up the present book. Finally, the author wishes to express his appreciation to Springer-Verlag for their most efficient handling of the publication of this book. 

Tokyo, September $1964$ 

Kôsaku Yosida 

\vspace{6mm}\begin{center}\bfseries Preface to the Second Edition\end{center}\vspace{2mm}
In the preparation of this edition, the author is indebted to Mr. \textsc{Floret} of Heidelberg who kindly did the task of enlarging the Index to make the book more useful. The errors in the second printing are corrected thanks to the remarks of many friends. In order to make the book more up-to-date, Section~\hyperref[sec:XIV.4]{$4$} of Chapter~\hyperref[chap:XIV]{XIV} has been rewritten entirely for this new edition. 

Tokyo, September $1967$ 

Kôsaku Yosida 

\vspace{6mm}\begin{center}\bfseries Preface to the Third Edition\end{center}\vspace{2mm}
A new Section ($9$. Abstract Potential Operators and Semi-groups) pertaining to G. \textsc{Hunt}'s theory of potentials is inserted in Chapter~\hyperref[chap:XIII]{XIII} of this edition. The errors in the second edition are corrected thanks to kind remarks of many friends, especially of Mr. \textsc{Klaus-Dieter Bier-Stedt}. 

Kyoto, April $1971$ 

Kôsaku Yosida 

\vspace{6mm}\begin{center}\bfseries Preface to the Fourth Edition\end{center}\vspace{2mm}
Two new Sections ``$6$. Non-linear Evolution Equations $1$ (The Kōmura-Kato Approach)'' and ``$7$. Non-linear Evolution Equations $2$ (The Approach Through The Crandall-Liggett Convergence Theorem)'' are added to the last Chapter~\hyperref[chap:XIV]{XIV} of this edition. The author is grateful to Professor Y. Kōmura for his careful reading of the manuscript. 

Tokyo, April $1974$ 

Kôsaku Yosida 

\vspace{6mm}\begin{center}\bfseries Preface to the Fifth Edition\end{center}\vspace{2mm}
Taking advantage of this opportunity, supplementary notes are added at the end of this new edition and additional references to books have been entered in the bibliography. The notes are divided into two categories. The first category comprises two topics: the one is concerned with the time reversibility of Markov processes with invariant measures, and the other is concerned with the uniqueness of the solution of time dependent linear evolution equations. The second category comprises those lists of recently published books dealing respectively with Sobolev Spaces, Trace Operators or Generalized Boundary Values, Distributions and Hyperfunctions, Contraction Operators in Hilbert Spaces, Choquet's Refinement of the Krein-Milman Theorem and Linear as well as Nonlinear Evolution Equations. 

A number of minor errors and a serious one on page $459$ in the fourth edition have been corrected. The author wishes to thank many friends who kindly brought these errors to his attention. 

Kamakura, August $1977$ 

Kôsaku Yosida 

\vspace{6mm}\begin{center}\bfseries Preface to the Sixth Edition\end{center}\vspace{2mm}
Two major changes are made to this edition. The first is the rewriting of the Chapter~\hyperref[chap:VI]{VI},$6$ to give a simplified presentation of Mikusiński's Operational Calculus in such a way that this presentation does not appeal to Titchmarsh's theorem. The second is the rewriting of the Lemma together with its Proof in the Chapter~\hyperref[chap:XII]{XII},$5$ concerning the Representation of Vector Lattices. This rewriting is motivated by a letter of Professor E. Coimbra of Universidad Nova de Lisboa kindly suggesting the author's careless phrasing in the above Lemma of the preceding edition. 

A number of misprints in the fifth edition have been corrected thanks to kind remarks of many friends. 

Kamakura, June $1980$ 

Kôsaku Yosida 

